%% =====================================================================
%% elision.tex -- replaces the theorem and proof sketch of Section 8 (paper-wo/main.tex, from
%% "\begin{theorem}[Elision soundness]" through the "\end{proof}" that follows it). Keeps label thm:elide;
%% Proof: appendix-proofs.tex, \S app:elide.
%% =====================================================================
\begin{theorem}[Elision soundness on the recorded execution]\label{thm:elide}
Let $P$ be a plan that removes or narrows barrier structs and keeps every layout transition, and let $\tau_P$ be the
trace with the plan applied. Let $F^*$ be the strengthened frame, in which each transition is split into a struct that
ends at the transition, the transition itself and a struct that starts from it, and load and store operations are
accesses. Let $F^*(\tau_P)$ satisfy the premises of Theorem~\ref{thm:sound}. If
$\Hpw(F^*(\tau_P)) \subseteq \Hpw(F^*(\tau))$, where pairs are compared by event identity (call, command and access
index), then $\Hspec(\tau_P) \subseteq \Hpw(F^*(\tau))$. In particular, if the pairwise mode reports nothing on
$F^*(\tau)$, then $\tau_P$ is race-free, and Corollary~\ref{cor:sched} applies to the transformed run on the recorded
host path.
\end{theorem}

The proof (Appendix~\ref{app:elide}) has two steps. First, $F^*$ only adds events and removes exemptions, and its
split structs chain exactly as the original struct did, so every race of a trace is a race of its strengthened frame.
Second, Theorem~\ref{thm:sound} on $F^*(\tau_P)$ and the premise give the inclusion. The premise is checked on the final
plan, so the order of the greedy search does not matter. Two requirements of the theorem are easy to lose in an
implementation. The comparison must be by event pairs. A comparison by diagnostic text or site can hide a new pair
behind an existing report of the same site, for example when a command buffer is submitted twice. It must also use the
pairwise mode. The epoch mode reports only a subset (Lemma~\ref{lem:modes}), so ``no new epoch report'' does not
imply the premise.

\paragraph{Beyond the recorded frames}
The theorem covers the frames that were traced, and nothing more. Reusing a plan on later frames would need the run to
continue as the recorded one, and a runtime matcher that falls back safely when it does not. We do not claim such a
reuse guarantee. Two obstacles show why a simple matcher is not enough: draining the device queues at a mismatch does
not drain the presentation engine, which may still read an image whose transition the plan narrowed; and a host write
can land in mapped memory before the matcher sees the call that reveals the mismatch. The layer
\code{VK\_LAYER\_COVENANT\_elide} matches each barrier call by its recorded key and applies the plan only to the
recorded ordinals; runs that go beyond them are measurements, not instances of Theorem~\ref{thm:elide}.
